\section{Introduction}
% A short explanation of how your method differs from, improves upon, or complements prior work. Begin with one or two sentences that articulate the motivations for the proposed methods. Then provide a concise description of the method’s design. Conclude with a summary of the major contributions (\eg, “Our contributions lies in three folds:” ).
% \begin{itemize}
%   \item First contribution.
%   \item Second contribution.
%   \item Third contribution.
% \end{itemize}

% introduction
While softmax attention has become a cornerstone of modern deep learning architectures, its $\mathcal{O}(n^2)$ computational complexity poses a significant bottleneck. This limitation is particularly pronounced in Vision Transformers (ViTs), where the number of visual tokens scales quadratically with the input image resolution.
% background
To address the issue, many works attempted to group vision tokens [CITE CSWin Transformer, Semantic equitable clustering, Dynamic group transformer]. Although this method reduces the complexity,
The Transformer loses its ability to perceive information globally.
Other works replace softmax attention with \textbf{linear attention} which helps reduce the quadratic complexity \cite{han2023flatten,fan2025breaking,gu2024mamba,cai2023efficientvit}.
However, linear attentions inability to model complex information yields unsatisfactory results due to its low-rank nature and magnitude negligence.
% transition
Many works have attempted to improve linear attention by tackling its low-rank nature \cite{fan2025breaking}.
Linear attention replaces softmax with kernel feature maps which causes the loss of sharpness.
Therefore, it can not be as selective as softmax attention which can place most of its attention weights on a few relevant tokens.
To address this, we apply a MARL gating mechanism to assign an importance score for each token, in which each agent will correspond to a visual token. We hypothesize that this will enable the agents to learn an optimal survival policy, ensuring that only the most salient and task relevant features are propagated into the recurrent memory. When trained from scratch and evaluated on the ImageNet100 dataset, the model achieved an accuracy of \(48.88\%\). Additionally, by transferring the trained backbone to CIFAR-10, we obtained an accuracy of \(58.37\%\). Although the proposed framework requires further refinement and the preliminary quantitative results remain modest, the generated saliency maps demonstrate that the gating mechanism successfully identifies important visual features.
% details


